%==============================================================================
\section{Biến thể và kỹ thuật nâng cao}
\label{sec:gd-part2}
%==============================================================================

\begin{ynghia}[title={Mở đầu}]
Nghiệm cuối cùng của GD phụ thuộc mạnh vào điểm khởi tạo và learning rate $\eta$. Ta trình bày các kỹ thuật thường dùng để khắc phục hạn chế đó --- \textbf{Momentum}, \textbf{Nesterov accelerated gradient (NAG)}, phân loại \textbf{Batch / SGD / mini-batch}, \textbf{điều kiện dừng} và ý tưởng \textbf{Newton's method}. Các biến thể này là công cụ chuẩn khi huấn luyện mô hình học máy, đặc biệt Deep Learning.
\end{ynghia}

%------------------------------------------------------------------------------
\subsection{Các thuật toán tối ưu Gradient Descent}
\label{subsec:gd-opt-algos}
%------------------------------------------------------------------------------

\subsubsection{Momentum}

\begin{phuongphap}[title={Nhắc lại thuật toán Gradient Descent}]
\begin{enumerate}
 \item Chọn điểm khởi tạo $\boldsymbol{\theta} = \boldsymbol{\theta}_0$.
 \item Lặp cập nhật cho đến khi đạt tiêu chí dừng:
 \[
 \boldsymbol{\theta} \leftarrow \boldsymbol{\theta} - \eta \nabla_{\boldsymbol{\theta}}J(\boldsymbol{\theta})
 \]
 với $\nabla_{\boldsymbol{\theta}}J(\boldsymbol{\theta})$ là gradient của hàm mất mát tại $\boldsymbol{\theta}$.
\end{enumerate}
\end{phuongphap}

\begin{ynghia}[title={Ẩn dụ vật lý: hòn bi trên bề mặt lồi}]
Thuật toán GD thường được minh họa bằng hình ảnh một hòn bi chịu trọng lực trên bề mặt dạng thung lũng. Bất kể đặt bi tại $A$ hay $B$, bi cuối cùng lăn xuống đáy $C$:

\mlcbimg{02_gradient_descent_part2/img03.png}
\end{ynghia}

\begin{luuy}[title={Hai đáy thung lũng và local minimum}]
Nếu bề mặt có hai đáy như hình (b), vị trí cuối cùng của bi phụ thuộc điểm khởi tạo: từ $A$ bi về $C$, từ $B$ bi về $D$. Điểm $D$ là \textbf{local minimum} không mong muốn khi $C$ là global minimum.
\end{luuy}

\begin{dongco}[title={Động lượng giúp vượt đáy cục bộ}]
Trong hình (b), nếu vận tốc ban đầu tại $B$ đủ lớn, bi có thể theo đà vượt qua $D$, leo dốc bên trái và (với động năng đủ lớn) tới $E$ rồi lăn xuống $C$ như hình (c). Ngược lại, bi lăn từ $A$ tới $C$ khó có đủ đà để vượt dốc $CE$ (dốc $CE$ cao hơn nhiều so với $DE$).

Hiện tượng này gợi ý thuật toán \textbf{Momentum} (theo đà): tích lũy hướng di chuyển qua các vòng lặp để giảm nguy cơ kẹt ở local minimum không mong muốn.
\end{dongco}

\begin{congthuc}[title={Cập nhật Gradient Descent với Momentum}]
Coi lượng thay đổi tại bước $t$ như vận tốc $v_t$; vị trí mới là $\boldsymbol{\theta}_{t+1} = \boldsymbol{\theta}_{t} - v_t$ (dấu trừ vì di chuyển ngược gradient). Vận tốc kết hợp độ dốc hiện tại và quán tính từ bước trước ($v_0 = \mathbf{0}$):
\[
 v_{t}= \gamma v_{t-1} + \eta \nabla_{\boldsymbol{\theta}}J(\boldsymbol{\theta}_t),
 \qquad
 \boldsymbol{\theta}_{t+1} = \boldsymbol{\theta}_t - v_t.
\]
Hệ số $\gamma$ thường chọn khoảng $0.9$; $\nabla_{\boldsymbol{\theta}}J(\boldsymbol{\theta}_t)$ là gradient tại $\boldsymbol{\theta}_t$.
\end{congthuc}

\begin{luuy}[title={Không gian nhiều chiều}]
Công thức trên áp dụng trực tiếp khi $\boldsymbol{\theta}$ là vector: các phép toán thực hiện theo từng tọa độ. Trong thực nghiệm, Momentum thường cải thiện đáng kể tốc độ và chất lượng hội tụ so với GD thuần.
\end{luuy}

\begin{vidu}[title={Hàm $f(x) = x^2 + 10\sin(x)$}]
Xét hàm một biến có hai local minimum, trong đó một điểm là global minimum:
\[
 f(x) = x^2 + 10\sin(x), \qquad f'(x) = 2x + 10\cos(x).
\]
Hàm này minh họa khả năng GD thuần dừng ở cực tiểu cục bộ trong khi Momentum có thể vượt qua.
\end{vidu}

\begin{vidu}[title={So sánh GD và GD với Momentum}]
Hình dưới thể hiện đường đi của nghiệm:

\mlcbframegrid{5}{02_gradient_descent_part2/img04}{0}{4}
\mlcbframegrid{5}{02_gradient_descent_part2/img05}{0}{100}

\textbf{Hàng trên} (không Momentum): hội tụ sau khoảng 5 vòng lặp nhưng dừng tại local minimum.

\textbf{Hàng dưới} (có Momentum): quỹ đạo vượt dốc tới vùng gần global minimum, dao động rồi giảm tốc và hội tụ. Số vòng lặp lớn hơn nhưng nghiệm chính xác hơn; đường đi phù hợp với mô hình vật lý có quán tính.
\end{vidu}

\begin{phuongphap}[title={Cài đặt \texttt{GD\_momentum} trong Python}]
Cho điểm khởi tạo \texttt{theta\_init}, hàm gradient \texttt{grad(theta)}, hệ số quán tính \texttt{gamma} và learning rate \texttt{eta}:
\begin{lstlisting}[language=Python]
# check convergence
def has_converged(theta_new, grad):
 return np.linalg.norm(grad(theta_new))/
 len(theta_new) < 1e-3

def GD_momentum(theta_init, grad, eta, gamma):
 # Suppose we want to store history of theta
 theta = [theta_init]
 v_old = np.zeros_like(theta_init)
 for it in range(100):
  v_new = gamma*v_old + eta*grad(theta[-1])
  theta_new = theta[-1] - v_new
  if has_converged(theta_new, grad):
   break
  theta.append(theta_new)
  v_old = v_new
 return theta
 # this variable includes all points in the path
 # if you just want the final answer,
 # use `return theta[-1]`
\end{lstlisting}
\end{phuongphap}

\begin{luuy}[title={Hạn chế của Momentum}]
\begin{itemize}
 \item \textbf{Không thoát được hố sâu:} nếu $\gamma$ quá nhỏ hoặc local minimum quá sâu/rộng, quán tính không đủ để vượt dốc.
 \item \textbf{Overshooting:} nếu $\gamma$ quá lớn, thuật toán có thể vượt qua global minimum, dao động mạnh và khó hội tụ chính xác tại đáy.
\end{itemize}
\end{luuy}

%------------------------------------------------------------------------------
\subsubsection{Nesterov accelerated gradient (NAG)}
%------------------------------------------------------------------------------

\begin{luuy}[title={Hạn chế của Momentum thuần}]
Momentum giúp vượt local minimum, nhưng gần nghiệm thuật toán thường dao động lâu trước khi dừng --- hệ quả trực tiếp của quán tính. \textbf{Nesterov accelerated gradient (NAG)} khắc phục bằng cách ``nhìn trước'' một bước, thường cho hội tụ nhanh hơn Momentum thuần.
\end{luuy}

\begin{ynghia}[title={Ý tưởng: gradient tại điểm nhìn trước}]
Thay vì dùng gradient tại $\boldsymbol{\theta}_t$, NAG ước lượng vị trí tiếp theo $\boldsymbol{\theta}_t - \gamma v_{t-1}$ (chỉ dùng thành phần momentum, chưa cộng gradient) rồi tính gradient tại điểm đó:

\mlcbimg{02_gradient_descent_part2/img06.jpeg}

\begin{itemize}
 \item \textbf{Momentum thường:} bước nhảy = momentum vector + gradient tại thời điểm hiện tại.
 \item \textbf{Nesterov:} bước nhảy = momentum vector + gradient tại điểm nhìn trước.
\end{itemize}
\end{ynghia}

\begin{congthuc}[title={Cập nhật NAG}]
\begin{align}
 v_{t} &= \gamma v_{t-1} + \eta \nabla_{\boldsymbol{\theta}}J\bigl(\boldsymbol{\theta}_t - \gamma v_{t-1}\bigr) \label{eq:nag-v}\\
 \boldsymbol{\theta}_{t+1} &= \boldsymbol{\theta}_t - v_{t} \label{eq:nag-theta}
\end{align}
Điểm đánh giá gradient khác Momentum thuần: $J$ được tính tại $\boldsymbol{\theta}_t - \gamma v_{t-1}$, không phải tại $\boldsymbol{\theta}_t$.
\end{congthuc}

\begin{vidu}[title={So sánh Momentum và NAG trên Linear Regression}]
\mlcbframegrid{5}{02_gradient_descent_part2/img07}{0}{81}
\mlcbframegrid{5}{02_gradient_descent_part2/img08}{0}{29}

\textbf{Hàng trên} (Momentum): đường đi zigzag, nhiều vòng lặp hơn.

\textbf{Hàng dưới} (NAG): hội tụ nhanh hơn, quỹ đạo ít zigzag hơn.
\end{vidu}

\begin{ynghia}[title={Adagrad, Adam, RMSprop --- tổng quan}]
Ngoài Momentum và NAG, Deep Learning thường dùng các optimizer \textbf{thích ứng learning rate theo từng tọa độ}: \textbf{Adagrad}, \textbf{RMSprop}, \textbf{Adam} (kết hợp momentum và RMSprop). Phần \textbf{Tổng hợp và củng cố} cuối giáo trình trình bày chi tiết hơn Adam và RMSprop.
\end{ynghia}

%------------------------------------------------------------------------------
\subsection{Biến thể của Gradient Descent}
\label{subsec:gd-variants}
%------------------------------------------------------------------------------

\begin{ynghia}[title={Bài toán Linear Regression (nhắc lại)}]
Hàm mất mát và gradient cho Linear Regression (ký hiệu $\mathbf{X}$ thay cho dữ liệu mở rộng $\bar{\mathbf{X}}$):
\[
 J(\mathbf{w}) = \frac{1}{2N}\|\mathbf{X}\mathbf{w} - \mathbf{y}\|_2^2
 = \frac{1}{2N} \sum_{i=1}^N(\mathbf{x}_i \mathbf{w} - y_i)^2,
 \qquad
 \nabla_{\mathbf{w}} J(\mathbf{w}) = \frac{1}{N}\sum_{i=1}^N \mathbf{x}_i^T(\mathbf{x}_i\mathbf{w} - y_i).
\]
\end{ynghia}

\subsubsection{Batch Gradient Descent}

\begin{dinhnghia}[title={Batch Gradient Descent}]
Thuật toán GD trong Chương 1 còn gọi là \textbf{Batch Gradient Descent}: mỗi bước cập nhật $\boldsymbol{\theta} = \mathbf{w}$ dùng \textbf{toàn bộ} $N$ điểm dữ liệu $\mathbf{x}_i$ để tính gradient.
\end{dinhnghia}

\begin{luuy}[title={Hạn chế với dữ liệu lớn và online learning}]
Với cơ sở dữ liệu rất lớn (ví dụ hàng tỉ bản ghi), tính lại gradient trên toàn bộ tập sau mỗi vòng lặp tốn kém và không phù hợp \textbf{online learning} --- mô hình cần cập nhật liên tục khi có dữ liệu mới. Batch GD buộc phải quét lại toàn bộ dữ liệu mỗi lần cập nhật, làm mất tính ``online''.

Giải pháp phổ biến hơn là \textbf{Stochastic Gradient Descent (SGD)}: mỗi bước chỉ dùng một (hoặc một nhóm nhỏ) mẫu.
\end{luuy}

\subsubsection{Stochastic Gradient Descent}

\begin{dinhnghia}[title={Stochastic Gradient Descent (SGD)}]
Tại mỗi bước, SGD tính gradient của hàm mất mát trên \textbf{một} điểm dữ liệu $\mathbf{x}_i$, cập nhật $\boldsymbol{\theta}$, rồi lặp cho các điểm khác. Một lượt duyệt hết dữ liệu gọi là một \textbf{epoch}. Thuật toán đơn giản nhưng hiệu quả trong thực tế.
\end{dinhnghia}

\begin{tinhchat}[title={Epoch và so sánh với Batch GD}]
\begin{itemize}
 \item Batch GD: mỗi epoch $\approx$ một lần cập nhật $\boldsymbol{\theta}$.
 \item SGD: mỗi epoch có $N$ lần cập nhật (mỗi mẫu một lần).
 \item Một epoch SGD có thể chậm hơn về thời gian duyệt, nhưng thường cần ít epoch hơn nhiều (ví dụ $\sim 10$ epoch ban đầu; với dữ liệu mới có thể dưới một epoch).
 \item SGD phù hợp dữ liệu lớn (Deep Learning) và bài toán online.
\end{itemize}
\end{tinhchat}

\begin{congthuc}[title={Quy tắc cập nhật SGD}]
\[
 \boldsymbol{\theta} \leftarrow \boldsymbol{\theta} - \eta \nabla_{\boldsymbol{\theta}} J(\boldsymbol{\theta}; \mathbf{x}_i; y_i),
\]
với $J(\boldsymbol{\theta}; \mathbf{x}_i; y_i)$ là mất mát trên một cặp $(\mathbf{x}_i, y_i)$. Sau mỗi epoch cần \textbf{xáo trộn} thứ tự mẫu để đảm bảo tính ngẫu nhiên; thứ tự ảnh hưởng đến hiệu năng.
\end{congthuc}

\begin{ynghia}[title={Kết hợp SGD với Momentum, AdaGrad, \ldots}]
Momentum, AdaGrad và các biến thể khác áp dụng trực tiếp lên SGD. Với Linear Regression ($\boldsymbol{\theta} = \mathbf{w}$):
\[
 J(\mathbf{w}; \mathbf{x}_i; y_i) = \frac{1}{2}(\mathbf{x}_i \mathbf{w} - y_i)^2,
 \qquad
 \nabla_{\mathbf{w}}J(\mathbf{w}; \mathbf{x}_i; y_i) = \mathbf{x}_i^T(\mathbf{x}_i \mathbf{w} - y_i).
\]
\end{ynghia}

\begin{phuongphap}[title={Cài đặt \texttt{sgrad} và \texttt{SGD} trong Python}]
\begin{lstlisting}[language=Python]
# single point gradient
def sgrad(w, i, rd_id):
 true_i = rd_id[i]
 xi = Xbar[true_i, :]
 yi = y[true_i]
 a = np.dot(xi, w) - yi
 return (xi*a).reshape(2, 1)

def SGD(w_init, grad, eta):
 w = [w_init]
 w_last_check = w_init
 iter_check_w = 10
 N = X.shape[0]
 count = 0
 for it in range(10):
 # shuffle data
 rd_id = np.random.permutation(N)
 for i in range(N):
  count += 1
  g = sgrad(w[-1], i, rd_id)
  w_new = w[-1] - eta*g
  w.append(w_new)
  if count%iter_check_w == 0:
   w_this_check = w_new
   if np.linalg.norm(w_this_check - w_last_check)/len(w_init) < 1e-3:
    return w
   w_last_check = w_this_check
 return w
\end{lstlisting}
\end{phuongphap}

\begin{vidu}[title={Kết quả SGD trên Linear Regression}]
Dữ liệu giống ví dụ Linear Regression ở Chương 1.

\textbf{Đường đi của nghiệm:}
\mlcbframegrid{5}{02_gradient_descent_part2/img09}{0}{144}

\textbf{Giá trị hàm mất mát trong 50 vòng lặp đầu:}
\mlcbimg{02_gradient_descent_part2/img10.png}

Đường đi zigzag hơn Batch GD vì mỗi bước chỉ phản ánh một mẫu; tuy vậy thuật toán nhanh chóng vào vùng lân cận nghiệm. Với $N = 1000$, SGD cần khoảng 3 epoch ($\sim 2911$ cập nhật), trong khi Batch GD tốt nhất trong thí nghiệm Chương 1 cần $\sim 49$ epoch --- lợi thế rõ rệt về tổng số lần cập nhật. Hàm mất mát trên toàn tập (hình cuối) giảm nhanh dù không mượt; chỉ với khoảng 10 mẫu đầu đã xấp xỉ được phương trình đường thẳng.
\end{vidu}

\subsubsection{Mini-batch Gradient Descent}

\begin{dinhnghia}[title={Mini-batch Gradient Descent}]
Mini-batch dùng $n$ mẫu mỗi bước ($1 < n \ll N$). Mỗi epoch: xáo trộn dữ liệu, chia thành các mini-batch kích thước $n$ (batch cuối có thể nhỏ hơn nếu $N$ không chia hết), rồi cập nhật lần lượt từng batch:
\[
 \boldsymbol{\theta} \leftarrow \boldsymbol{\theta} - \eta\nabla_{\boldsymbol{\theta}} J(\boldsymbol{\theta}; \mathbf{x}_{i:i+n}; \mathbf{y}_{i:i+n}),
\]
với $\mathbf{x}_{i:i+n}$ là $n$ mẫu liên tiếp sau khi xáo trộn (ký hiệu kiểu Python).
\end{dinhnghia}

\begin{luuy}[title={Kết hợp với Momentum, Adagrad, Adadelta, \ldots}]
Các kỹ thuật tăng tốc GD (Momentum, Adagrad, Adadelta, \ldots) áp dụng tương tự trên mini-batch.
\end{luuy}

\begin{ynghia}[title={Ứng dụng thực tế}]
Mini-batch GD là chuẩn trong Machine Learning, đặc biệt Deep Learning. Kích thước batch $n$ thường chọn trong khoảng $50$--$100$: cân bằng độ ổn định gradient (gần Batch) và chi phí tính toán (gần SGD).
\end{ynghia}

\begin{vidu}[title={Dao động hàm mất mát}]
Giá trị mất mát sau mỗi lần cập nhật $\boldsymbol{\theta}$ trong một bài toán phức tạp hơn:

\mlcbimg{02_gradient_descent_part2/img11.png}

Hàm mất mát dao động qua từng bước nhưng có xu hướng giảm và hội tụ. Tài liệu tham khảo chi tiết: \href{http://sebastianruder.com/optimizing-gradient-descent/index.html#stochasticgradientdescent}{Stochastic Gradient Descent} (Sebastian Ruder).
\end{vidu}

%------------------------------------------------------------------------------
\subsection{Điều kiện dừng}
\label{subsec:stopping-criteria}
%------------------------------------------------------------------------------

\begin{ynghia}[title={Khi nào dừng thuật toán?}]
Tiêu chí dừng quyết định thời điểm kết thúc lặp. Trong thực nghiệm thường dùng một hoặc kết hợp các cách sau:

\begin{enumerate}
 \item \textbf{Giới hạn số vòng lặp:} phổ biến nhất, đảm bảo thời gian chạy có hạn; nhược điểm là có thể dừng trước khi đủ gần nghiệm.
 \item \textbf{Chuẩn gradient:} dừng khi $\|\nabla_{\boldsymbol{\theta}} J\|$ đủ nhỏ. Với SGD/mini-batch, tính gradient toàn tập tốn kém nên ít dùng trực tiếp.
 \item \textbf{Thay đổi hàm mất mát:} dừng khi $|J(\boldsymbol{\theta}_{t+1}) - J(\boldsymbol{\theta}_t)|$ nhỏ. Có thể dừng sớm tại \textbf{saddle point} (vùng phẳng cục bộ không phải cực tiểu).
 \item \textbf{So sánh nghiệm (SGD, mini-batch):} kiểm tra $\|\boldsymbol{\theta}_{t+k} - \boldsymbol{\theta}_t\|$ sau mỗi $k$ bước (ví dụ $k = 10$ trong mã SGD phía trên) --- thường hiệu quả trong thực hành.
\end{enumerate}
\end{ynghia}

%------------------------------------------------------------------------------
\subsection{Newton's method}
\label{subsec:newton}
%------------------------------------------------------------------------------

\begin{ynghia}[title={Phương pháp bậc nhất và bậc hai}]
Các phương pháp GD đã trình bày là \textbf{first-order methods} (chỉ dùng gradient bậc một). \textbf{Newton's method} là \textbf{second-order method}: bước cập nhật còn dùng đạo hàm bậc hai (Hessian).

Nhắc lại: cực tiểu thỏa $\nabla J = \mathbf{0}$. Newton's method nổi tiếng trong việc giải phương trình $f(x) = 0$; áp dụng cho $f'(x) = 0$ ta có phương pháp tối ưu Newton.
\end{ynghia}

\subsubsection{Newton's method cho $f(x) = 0$}

\begin{ynghia}[title={Ý tưởng hình học}]
Từ $x_0$ gần nghiệm $x^*$, vẽ tiếp tuyến $y = f(x)$ tại $x_t$; giao điểm $x_{t+1}$ với trục hoành được coi là xấp xỉ tốt hơn. Lặp đến khi $f(x_t) \approx 0$:

\mlcbimg{02_gradient_descent_part2/img12.png}
\end{ynghia}

\begin{congthuc}[title={Công thức Newton một biến}]
Tiếp tuyến tại $x_t$: $y = f'(x_t)(x - x_t) + f(x_t)$. Giao với trục $x$ (đặt $y = 0$):
\[
 x_{t+1} = x_t - \frac{f(x_t)}{f'(x_t)}.
\]
\end{congthuc}

\subsubsection{Newton cho tìm cực tiểu}

\begin{congthuc}[title={Áp dụng cho $f'(x) = 0$ và không gian nhiều chiều}]
Giải $f'(x) = 0$ bằng Newton:
\[
 x_{t+1} = x_t -(f''(x_t))^{-1}{f'(x_t)}.
\]
Với vector tham số $\boldsymbol{\theta}$:
\[
 \boldsymbol{\theta}_{t+1} = \boldsymbol{\theta}_t - \mathbf{H}(J(\boldsymbol{\theta}_t))^{-1} \nabla_{\boldsymbol{\theta}} J(\boldsymbol{\theta}_t),
\]
trong đó $\mathbf{H}(J)$ là \textbf{ma trận Hessian} (đạo hàm bậc hai của $J$).
\end{congthuc}

\subsubsection{Hạn chế của Newton's method}

\begin{luuy}[title={Ba hạn chế chính}]
\begin{enumerate}
 \item \textbf{Điểm khởi tạo phải gần nghiệm.} Khai triển Taylor bậc một: $0 = f(x^*) \approx f(x_t) + f'(x_t)(x_t - x^*)$, suy ra $x^* \approx x_t - f(x_t)/f'(x_t)$ chỉ đúng khi $x_t$ gần $x^*$. Có ví dụ phân kỳ kinh điển: nghiệm gần $-2$, nhưng khởi tạo tại $0$ và $1$ luân phiên, không hội tụ.

 \item \textbf{Đạo hàm mẫu số gần $0$.} Khi $f'(x_t) \approx 0$, tiếp tuyến gần song song trục $x$; giao điểm ở vô cực hoặc không xác định. Đặc biệt nghiêm trọng nếu nghiệm nằm tại điểm có đạo hàm bằng $0$.

 \item \textbf{Chi phí Hessian trong không gian lớn.} Tính và nghịch đảo $\mathbf{H}$ tốn bộ nhớ và thời gian khi số chiều và số mẫu lớn --- hạn chế lớn so với GD bậc một trong Deep Learning.
\end{enumerate}
\end{luuy}

%------------------------------------------------------------------------------
\subsection{Tổng kết}
\label{subsec:gd-part2-conclusion}
%------------------------------------------------------------------------------

\begin{tongket}[title={Tổng kết chương}]
Chương này bổ sung \textbf{Momentum}, \textbf{NAG}, phân loại \textbf{Batch / SGD / mini-batch}, \textbf{điều kiện dừng} và ý tưởng \textbf{Newton}. Đây là các công cụ thực tế khi huấn luyện mô hình quy mô lớn. Phần \textbf{Tổng hợp và củng cố} (cuối giáo trình) gom bảng so sánh, Adam/RMSprop và câu hỏi ôn tập.
\end{tongket}

\begin{phuongphap}[title={Tài liệu tham khảo}]
\begin{enumerate}
 \item Newton's method --- Wikipedia: \url{https://en.wikipedia.org/wiki/Newton%27s_method}
 \item An overview of gradient descent optimization algorithms: \url{http://sebastianruder.com/optimizing-gradient-descent/}
 \item Stochastic Gradient Descent --- Wikipedia: \url{https://en.wikipedia.org/wiki/Stochastic_gradient_descent}
 \item Stochastic Gradient Descent --- Andrew Ng (Coursera lecture notes)
 \item Nesterov, Y. (1983). A method for unconstrained convex minimization problem with the rate of convergence $O(1/k^2)$. \emph{Doklady ANSSSR} (translated as \emph{Soviet Math. Dokl.}), vol.~269, pp.~543--547.
\end{enumerate}
\end{phuongphap}
